\documentclass[11pt]{article}

\usepackage[margin=1in]{geometry}
\usepackage{amsmath,amssymb,amsthm}
\usepackage{booktabs}
\usepackage{graphicx}
\usepackage{xcolor}
\usepackage{multirow}
\usepackage{hyperref}
\usepackage{natbib}
\usepackage{setspace}
\doublespacing
\usepackage{lineno}
\linenumbers

\newtheorem{proposition}{Proposition}

\renewcommand{\arraystretch}{1.3}

\title{Supplementary Materials\\[6pt]
\large Federated site-level pooling estimates a different quantity
than patient-level analysis: ecological bias across
42 million COVID-19 records}

\author{Elliot Tower\\
\textit{Independent Researcher}\\
\texttt{elliot@elliottower.ai}}

\date{}

\begin{document}
\maketitle

\setcounter{section}{0}
\renewcommand{\thesection}{Supplement~\Alph{section}}


\section{Coverage-Degeneracy of the Consistency Test}
\label{supp:degeneracy}

The consistency statistic becomes degenerate (null variance $= 0$)
when all sites report all variable pairs (homogeneous coverage).
The reason is that the within-dimension permutation scheme shuffles
correlations independently for each variable pair, and when every
site pair shares the same set of variable pairs, each per-dimension
contribution to the statistic is invariant under permutation.
Heterogeneous coverage breaks this invariance because permutation
changes which site-pair sums include which terms.

This is a general property of permutation tests: when the data
structure has a symmetry that the permutation scheme preserves, the
null distribution collapses~\cite{phipson2010permutation}.

We confirmed this empirically on all three datasets.  Both the CDC
pseudo-sites (9 groups, all 15 variable pairs) and Mexico's 32 states
(all 15 pairs) have homogeneous coverage, producing degenerate
permutation nulls ($z \approx 0$, null standard deviation
$< 10^{-16}$).  The 4CE data has heterogeneous coverage (10 sites
report 36--296 of 1,305 possible comorbidity pairs; 10 report all
1,305), producing $z = 25$ ($p < 0.0001$).  The degeneracy on
Mexico's data---which uses real geographic sites, not constructed
pseudo-sites---confirms that the limitation is structural (homogeneous
coverage), not an artifact of the CDC's pseudo-site design.
Supplement~S10 provides additional calibration: under the null,
2{,}000 simulations with heterogeneous coverage produce a
false-positive rate of 4.1\% at $|z| > 2$, confirming the test is
well-calibrated when coverage conditions are met.


\section{Pre-Registered Supplementary Analyses}
\label{supp:preregistered}

The following ten analyses were specified in a frozen protocol
(commit \texttt{97f7094}) before any results were examined.  Each
specifies a falsification criterion; two fired (\dag).
Table~\ref{tab:supplement_summary} summarizes all results.

\begin{table}[h]
\centering
\caption{Summary of pre-registered supplementary analyses.
\dag~indicates a fired falsification criterion.  All results are
reported regardless of outcome.}
\label{tab:supplement_summary}
\small
\begin{tabular}{clll}
\toprule
\# & Analysis & Key result & Falsified? \\
\midrule
S1 & Power simulation (CDC) & Power = 100\% all conditions & No \\
S2\dag & Mundlak decomposition (Mexico) & Within OR = 5.22; between unidentified & Yes \\
S3 & Expected ecological (Mexico) & 2.09$\times$ amplification (109\% bias) & No \\
S4 & Duncan--Davis bounds (CDC) & Bounds [0, 24.4\%], contains truth & No \\
S5\dag & Positive control (CDC) & Sex also fails ($p = 0.99$) & Yes \\
S6 & Alternative pseudo-sites (CDC) & Random 9-group power = 4.8\% & No \\
S7 & E-value sensitivity (4CE) & E-value $= 5.8 \times 10^{5}$ & No \\
S8 & Leave-one-out stability (4CE) & RE/FE ratio 18--984 after LOO & No \\
S9 & Two-stage vs.\ one-stage (Mexico) & OR = 5.34 vs.\ 5.22 (2\% agreement) & No \\
S10 & Consistency calibration & FPR = 4.1\% at $|z|>2$ & No \\
\bottomrule
\end{tabular}
\end{table}

\paragraph{S1: Power vs.\ bias decomposition (CDC).}
We simulated individual-level data within each of the 7 CDC sites using
known ORs (age $= 9.88$, comorbidity $= 1.5$, sex $= 1.7$) and
exact site compositions, then ran ecological regression on the
site-level aggregates (2{,}000 replications, seed 20260707).
Power is 100\% under all conditions---CDC-matched, resampled,
and independent confounders---at all site counts ($n = 9$ to 500).
\emph{Falsification:} the simulation confirms the ecological
regression should detect the effect, consistent with the observed
significance ($p = 0.02$).

\paragraph{S2: Mundlak within/between decomposition (Mexico).}
A logistic regression with state fixed effects and within-site/between-site
elderly covariates on 229{,}125 Mexico records yields a within-site
coefficient of $+1.65$ (OR $= 5.22$, 95\% CI [4.99, 5.47]).  The
between-site coefficient is unidentified (SE $\to \infty$) because the
state means are collinear with the 31 state dummies.
\emph{Falsification:} the CI-overlap criterion fires, but the collinearity
is a specification property, not evidence against aggregation bias.
The within-site OR of 5.22 matches S9's one-stage estimate exactly.

\paragraph{S3: Expected ecological coefficient (Mexico).}
The individual-level logistic model (elderly OR $= 3.25$, male
OR $= 0.61$, comorbidity OR $= 3.00$) predicts an ecological slope
of $+0.11$ by integrating over each state's covariate distribution.
The observed ecological slope is $+0.23$---2.09 times the expectation,
a 109\% amplification from aggregation bias.

\paragraph{S4: Duncan--Davis ecological-inference bounds (CDC).}
Across 6 valid CDC sites (excluding those with zero elderly proportion),
the intersection of Duncan--Davis bounds on the elderly mortality rate
is [0\%, 24.4\%].  The observed rate (9.91\%) falls within these bounds,
but the 24.4-percentage-point width confirms that ecological data alone
cannot pin down the elderly-specific mortality rate.

\paragraph{S5: Reduced-confounding positive control (CDC).}
Sex ratio---a predictor with simpler confounding
($\rho(\text{elderly}, \text{male}) = +0.41$,
$\rho(\text{elderly}, \text{comorbidity}) = -0.05$)---also fails
to predict mortality ecologically on CDC data (slope $= +0.0009$,
$p = 0.99$).  \emph{Falsification:} even a predictor with low
confounding fails, indicating the pseudo-site construction is too
noisy for any ecological regression at $n = 9$.

\paragraph{S6: Alternative pseudo-site definitions (CDC).}
We randomly partitioned 105 million CDC records into 9 groups
(500 replications, seed 20260708) and ran ecological regression
on each partition.  Power is 4.8\% (95\% Wilson CI [3.2\%, 7.0\%])---at
the false-positive rate.  With 50 groups, power rises to 29.2\%
(CI [25.4\%, 33.3\%]).  The CDC null is a general property of
ecological regression at small $k$.

\paragraph{S7: E-value sensitivity analysis (4CE).}
The E-value for the elderly-proportion moderator ($\beta = 14.2$,
SE $= 0.83$) has a point estimate of $2.9 \times 10^{6}$ and a
95\% CI lower bound of $5.8 \times 10^{5}$
\cite{vanderweele2017}.  No plausible unmeasured confounder could
explain this association.

\paragraph{S8: Leave-one-site-out stability (4CE).}
Across 20 sites at 90 days, the full-dataset RE/FE ratio for CNS PMI
is 534:1.  Dropping each site in turn, the ratio ranges from 18.2
(dropping ICSM) to 983.7 (dropping NWU), never falling below 10-fold.
No single site drives the divergence.

\paragraph{S9: Two-stage vs.\ one-stage IPD (Mexico).}
Three estimators on Mexico's 32 states: (a) ecological regression
($\beta = +0.23$); (b) two-stage IPD (per-state logistic pooled
via RE meta-analysis: OR $= 5.34$, 95\% CI [4.66, 6.13],
$I^2 = 0.88$); (c) one-stage IPD (logistic with state fixed effects:
OR $= 5.22$, 95\% CI [4.98, 5.46]).  The two IPD methods agree
within 2\%.  Ecological regression produces an incommensurable quantity.

\paragraph{S10: Consistency-test calibration.}
Under the null (no cross-site heterogeneity), 2{,}000 simulations
with 30 binary variables, 435 pairwise correlations, and 20 sites
(10 full-coverage, 10 partial) produce a false-positive rate of
4.1\% at $|z| > 2$ (target $\leq 6\%$) and 0\% at $|z| > 25$.
The consistency test is well-calibrated.


\paragraph{Table S1: Cross-dataset method comparison.}

\begin{table}[h]
\centering
\caption*{\textbf{Table S1.} Summary of methods applied across all three datasets.  The
consistency statistic has power only on heterogeneous-coverage data (4CE).
Per-site interaction analysis detects effect modification on all three
datasets.  Individual-level analysis (where available) gives the true
answer.}
\label{tab:comparison}
\begin{tabular}{llcc}
\toprule
Method & Dataset & Statistic & $p$ \\
\midrule
\multirow{3}{*}{Ecological regression}
  & CDC (7 groups) & $\beta = +0.55$ & 0.02 \\
  & Mexico (32 states) & $\beta = +1.31$ & $< 0.0001$ \\
  & 4CE (21 hospitals) & PMI $= 2.0$ & $< 10^{-6}$ \\
\midrule
\multirow{3}{*}{Consistency statistic}
  & CDC & degenerate & --- \\
  & Mexico & degenerate & --- \\
  & 4CE & $z = 25$ & $< 0.0001$ \\
\midrule
\multirow{3}{*}{Per-site interaction}
  & CDC & OR: 0.23--0.69 & all $< 10^{-4}$ \\
  & Mexico & OR: 0.19--0.38 & all $< 10^{-4}$ \\
  & 4CE & $\beta$: $+32$ to $+38$ & $< 0.04$ \\
\midrule
\multirow{2}{*}{Individual logistic}
  & CDC & OR $= 9.88$ & $< 0.001$ \\
  & Mexico & OR $= 11.26$ & $< 0.001$ \\
\bottomrule
\end{tabular}
\end{table}


\bibliographystyle{unsrtnat}
\bibliography{references}

\end{document}
